\documentclass{article}
\usepackage{amsmath, amssymb, amsthm}
\usepackage{geometry}
\usepackage[UTF8]{ctex} % 用于中文支持，如果不需要中文，可以移除此行

\newtheorem{problem}{问题}
\newtheorem{solution}{解答}
\title{数值分析第一周作业}
\date{\today}
\author{李佩优PB23000270}
\geometry{a4paper,scale=0.8}
\begin{document}
\maketitle
\section*{2}
对于给定的函数$f(x)$，在n+1个指定的结点上次数$\leqslant n$的插值多项式$p$是唯一确定的。
因此，存在一个映射$f\mapsto p$。用L标势这个映射，并证明
$$Lf = \sum_{i=0}^{n} f(x_i) \ell_i$$
再证明L是现行的，即$L(af+bg)=aLf+bLg$。
\begin{proof}

注意到拉格朗日基函数：$\ell_i(x)=\prod_{\substack{0\leq j\leq n\\ j\neq i}} \frac{x-x_j}{x_i-x_j}$满足：
$j\neq i$时，$\ell_i(x_j)=0$; $\ell_i(x_i)=1$，且为n次多项式。
于是，多项式$$p(x)=\sum_{i=0}^{n} f(x_i) \ell_i(x)$$
在结点$x_i$处满足$$p(x_i)=\sum_{i=0}^{n} f(x_i) \ell_i(x_i)=f(x_i)$$，且其次数不超过n。
因此$p$是满足条件的插值多项式，由插值多项式的唯一性可知$p$就是$Lf$，即$$Lf = \sum_{i=0}^{n} f(x_i) \ell_i.$$

接下来证明L是线性的。
对于任意的函数$f$和$g$，以及任意的常数$a$和$b$，
考虑$L(af+bg)$,$\forall i, aLf(x_i)+bLg(x_i)=af(x_i)+bg(x_i)=(af+bg)(x_i)$，
$\Rightarrow aLf+bLg$也是$af+bg$的插值多项式。由于插值多项式的唯一性可知$L(af+bg)=aLf+bLg$。

\end{proof}
\section*{3}
G: $$Gf = \sum_{i=0}^{n}f(x_i)\ell_i^2$$
\begin{proof}
由$\ell_i$是n次多项式可知$\ell_i^2$是2n次多项式，因此$Gf$是次数不超过2n次的多项式。

$\forall i\neq j, \ell_i(x_j)=0;\quad \ell_i(x_i)=1$;
$\Rightarrow \forall i\neq j, \ell_i^2(x_j)=0;\quad \ell_i^2(x_i)=1$。
$$\Rightarrow Gf(x_i)=\sum_{i=0}^{n}f(x_i)\ell_i^2(x_i)=f(x_i)$$
因此$Gf$也是满足条件的插值多项式。

f非负时，$\forall i, f(x_i)\geq 0$，
$$\Rightarrow Gf(x)=\sum_{i=0}^{n}f(x_i)\ell_i^2(x)\geq 0$$
因此$Gf$也是非负的。

\end{proof}
\section*{4}
对每个次数最多是n次的多项式$q$，都有$Lq=q$。
\begin{proof}
假设$Lq\neq q$
，那么$Lq-q$是一个次数不超过n次的多项式，并且在结点$x_i$处满足$(Lq-q)(x_i)=0$。
因此$Lq-q$有n+1个不同的根，但其次数不超过n，这与多项式的根的个数不超过其次数的定理矛盾。
因此$Lq=q$。
\end{proof}
\section*{5}
证明对所有x,都有$$\sum_{i=0}^{n} \ell_i(x)=1$$
\begin{proof}
$\forall i\neq j, \ell_i(x_j)=0;\quad \ell_i(x_i)=1$;
$\Rightarrow \forall i, \sum_{i=0}^{n} \ell_i(x_i)=1$。
因此常函数1本身是一个次数不超过n的多项式，因此1是它的插值多项式。
又由$\sum_{i=0}^{n} \ell_i(x)$是次数不超过n的多项式，
$\Rightarrow 1=L(\sum_{i=0}^{n} \ell_i(x))=\sum_{i=0}^{n} \ell_i(x)$。
即$\sum_{i=0}^{n} \ell_i(x)\equiv 1$。
\end{proof}
\section*{6}
$\forall x ,\sum_{i=0}^{n} \ell_i(x)=1$,
$\Rightarrow \sum_{i=0}^{n} f(x)\ell_i(x) = f(x)$。
又有$p=Lf = \sum_{i=0}^{n} f(x_i) \ell_i$
$\Rightarrow f(x)-p(x)=\sum_{i=0}^{n} (f(x)-f(x_i))\ell_i(x)$。
\section*{8}
对于给定的$p(0),p(1)$和$p'(\xi )$,其中$\xi$是任意预先给定的点。
讨论求一个次数至多为2次的多项式问题。

\begin{solution}
令 \(p(x)=ax^2+bx+c\)，则由条件得
\[
c=\alpha,\quad a+b=\beta-\alpha,\quad 2a\xi+b=\gamma.
\]
解此方程组：

- 当 \(\xi\neq\frac{1}{2}\) 时，有唯一解
  \[
  a=\frac{\gamma-(\beta-\alpha)}{2\xi-1},\quad b=\beta-\alpha-a,\quad c=\alpha,
  \]
  即
  \[
  p(x)=\frac{\gamma-(\beta-\alpha)}{2\xi-1}x^2+\left(\beta-\alpha-\frac{\gamma-(\beta-\alpha)}{2\xi-1}\right)x+\alpha.
  \]

- 当 \(\xi=\frac{1}{2}\) 时，方程组有解当且仅当 \(\gamma=\beta-\alpha\)。此时 \(a\) 可取任意常数，\(b=\beta-\alpha-a\)，\(c=\alpha\)，即解集为
  \[
  p(x)=ax^2+(\beta-\alpha-a)x+\alpha,\quad \forall a\in\mathbb{R}\ (\text{或 }\mathbb{C}).
  \]
  若 \(\gamma\neq\beta-\alpha\)，则无解。
\end{solution}

\end{document}